% Generated draft: requires mathematical and source review.
\documentclass[11pt]{article}
\usepackage[margin=25mm]{geometry}
\usepackage{fontspec}
\setmainfont{Latin Modern Roman}
\usepackage{amsmath,amssymb,mathtools,unicode-math}
\setmathfont{Latin Modern Math}
\usepackage{enumitem,needspace,xurl,hyperref}
\setlength{\parindent}{0pt}
\setlength{\parskip}{5pt}
\setlength{\emergencystretch}{3em}
\newcommand{\N}{\mathbb{N}}
\newcommand{\Z}{\mathbb{Z}}
\newcommand{\Q}{\mathbb{Q}}
\newcommand{\R}{\mathbb{R}}
\newcommand{\C}{\mathbb{C}}
\newcommand{\F}{\mathbb{F}}
\newcommand{\A}{\mathbb{A}}
\newcommand{\PP}{\mathbb{P}}
\newcommand{\E}{\mathbb{E}}
\newcommand{\eps}{\varepsilon}
\newcommand{\dd}{\,\mathrm{d}}
\begin{document}
\section*{UnsolvedMath ID: 30003234}
\textbf{Source problem code:} catalogue entry
\subsection{Definitions and mathematical statement}
For \(t\in\R\), write \(\|t\|=\min_{p\in\Z}|t-p|\). To make the notation explicit, provisionally interpret \(\operatorname{Bad}(i,j)\) using the standard simultaneous weighted approximation convention. For \(i,j\geq0\) with \(i+j=1\), put \(w_1=i\), \(w_2=j\), and define
\[
\operatorname{Bad}(i,j)=\left\{(x_1,x_2)\in\R^2:\ \exists c>0\ \forall q\in\Z_{\geq1},\quad \max_{\substack{k\in\{1,2\}\\w_k>0}}q^{w_k}\|qx_k\|\geq c\right\}.
\]
Thus a coordinate of weight zero imposes no condition. In particular, if
\[
\operatorname{Bad}_1=\{x\in\R:\exists c>0\ \forall q\in\Z_{\geq1},\ q\|qx\|\geq c\},
\]
then \(\operatorname{Bad}(0,1)=\R\times\operatorname{Bad}_1\) and \(\operatorname{Bad}(1,0)=\operatorname{Bad}_1\times\R\). The supplied record does not define the weighted set or its endpoint conventions, so this interpretation requires confirmation against the source.

The supplied question is: given a planar curve \(C\subseteq\R^2\) and weights \(i,j\geq0\) satisfying \(i+j=1\), is
\[
\bigl|C\cap\operatorname{Bad}(i,j)\bigr|>\aleph_0?
\]
The intended class of planar curves is unspecified. Consequently, a precise universal formulation requires additional hypotheses on \(C\); the question above records the supplied target without asserting an affirmative answer for every planar curve. No regularity, non-degeneracy, or restrictions on lines may silently be added.

The supplied dated curation describes the question as partially solved: it reports full curve dimension for the intersection on twice-differentiable non-degenerate planar curves, results for several lines, and unresolved or empty cases among arbitrary planar curves. These are supplied literature qualifications, not independently verified mathematical assertions; the exact hypotheses and meaning of non-degeneracy are not supplied.
\subsection{Short English statement}
For a planar curve and nonnegative weights summing to one, the question asks whether uncountably many points on the curve are weighted badly approximable. The intended curve class and the source's definition at zero weights require clarification. The supplied partially solved status and reported results are dated curation requiring independent review.
\subsection{Sources}
\begin{itemize}
\item Supplied local UnsolvedMath registry record, permanent ID \(30003234\), title: Badly Approximable Points on Planar Curves; collection: Oberwolfach Reports Open Problems; exact source problem code: \texttt{catalogue entry}. The project registry locations are \texttt{data/unsolvedmath/problems.json} and \texttt{catalog/problems.json}; this draft uses the supplied complete record.
\item Supplied report citation: Arbeitsgemeinschaft: Diophantine Approximation, Fractal Geometry and Dynamics (2017), Oberwolfach Reports, pp.\ 2749--2792, DOI \texttt{10.4171/owr/2016/48}. \url{https://doi.org/10.4171/owr/2016/48}. The registry records proposed year 2016. No individual proposer is supplied; the background calls this Davenport's question.
\item Supplied literature reference: Badziahin--Velani, badly approximable points on planar curves. \url{https://doi.org/10.1007/s00208-014-1020-z}.
\item Supplied literature reference: Oberwolfach Report 48. \url{https://doi.org/10.4171/owr/2016/48}.
\item Literature checked at: \texttt{2026-08-21}. This exact date refers to supplied dated curation, not independent verification. Its assessment labels the problem partially solved and reports results for twice-differentiable non-degenerate planar curves with full curve dimension, and for several lines, while arbitrary planar curves include unresolved or empty cases. The background explicitly requests independent verification of time-sensitive status and cited claims.
\item The registry labels the statement \texttt{corrected\_verified} and reports manual semantic, typographic, and LaTeX review. These supplied labels do not constitute independent verification of the formulation, cited results, or mathematical correctness of this draft.
\end{itemize}
\end{document}
